\head{\textsc{\emph{Ilíada}, Canto II}}
\section*{\centering\textbf{\textsc{Canto} II}\footnote{\neh, pp. 92-93.}}
\addcontentsline{toc}{subsection}{\textsc{Canto II}}
%
\noindent-- -- -- -- -- --\par
%
\beginnumbering
\pstart
\setline{474}Tal como los cabreros distinguen fácilmente\\
las copiosas majadas, si en el pasto se mezclan,\\
así les dan los jefes el orden conveniente\\
para el combate. Entre ellos va el rey Agamenón,\\
parecido por ojos y testa al dios tonante;\\
por el talle a Ares y por el pecho a Poseidón.\\
Cual la mayor cabeza del hato es el pujante\\
toro que entre las vacas sobresale, aquel día\\
dispuso Zeus que Atrida fuera el más arrogante\\
de porte entre los muchos héroes que allí había.\par
\pend
\endnumbering
%
\noindent-- -- -- -- -- --\par